\section{Conclusion}
\label{sec:conclusion}

We presented \Weld, an integrated developer workspace with a built-in study
mode, and a hybrid within/between-subjects user study ($N$~=~\emph{TBD}) that
isolates the effect of agent context breadth on task performance, cognitive
load, user reliance, and the calibration of trust against measured
effectiveness. The study's primary methodological
contribution is the \ContextGate: a server-side, code-enforced context-scoping
mechanism that makes the agent's information diet both configurable and
auditable, enabling a clean between-subjects manipulation without relying on
prompt engineering.

% TODO: Fill in empirical conclusions once data collection is complete.

Our results [to be reported] suggest [summary of key empirical findings].
Together, these findings have implications for the design of AI agents in
collaborative productivity tools, particularly regarding (1) the value of
gradually expanding context rather than granting full information access by
default, and (2) the importance of audit-logging context assembly decisions as a
trust and governance mechanism.

Future work should examine longitudinal effects of integrated AI agents on team
coordination patterns, study generalization to professional teams working on
real projects, and explore adaptive context scoping---where the agent's
information diet expands or contracts based on task characteristics detected at
runtime.

\Weld and the study infrastructure are open-source and available at
\url{https://github.com/weldtool/weld}.
